% =======================================================
% 3. MODELLO DI SVILUPPO
% =======================================================
\section{Modello di Sviluppo}

\subsection{Scelta del modello}
Il gruppo adotta un modello di sviluppo agile ispirato a Scrum, con
sprint settimanali da martedì a martedì e rotazione settimanale dei
ruoli. La scelta è motivata dalla natura incrementale del progetto e dalla necessità di
adattarsi rapidamente ai feedback del committente. 

\subsection{Motivazioni}
\begin{itemize}
    \item \textbf{Adattabilità:} il piano di lavoro può essere revisionato sprint per
          sprint in base ai risultati e ai feedback ottenuti.
    \item \textbf{Visibilità continua:} le cerimonie strutturate garantiscono che tutto
          il team abbia sempre un quadro aggiornato dello stato di avanzamento.
    \item \textbf{Qualità incrementale:} la Definition of Done (Sez.~\ref{sec:dod}) e
          le code review sistematiche riducono il rischio di accumulare debito tecnico.
    \item \textbf{Sviluppo competenze trasversali:} la rotazione settimanale dei ruoli
          garantisce che ogni membro acquisisca esperienza su tutti gli aspetti del
          progetto.
\end{itemize}

\subsection{Cerimonie e cadenza settimanale}
Ogni sprint ha durata di una settimana (martedì → martedì). La riunione principale si
tiene il \textbf{martedì alle 14:30} e segue un ordine fisso:
\begin{enumerate}
    \item \textbf{Sprint Retrospective} dello sprint precedente: cosa ha funzionato,
          cosa no, azioni correttive da applicare subito.
    \item \textbf{Sprint Review} dello sprint precedente: verifica di ciò che è stato
          prodotto rispetto alla Definition of Done, chiusura formale dei task.
    \item \textbf{Sprint Planning}: aggiornamento del backlog prioritario, stima del
          tempo previsto, assegnazione dei task tramite Google Form.
          Output: Sprint Backlog della settimana.
\end{enumerate}

\begin{longtable}{p{2.5cm} p{2cm} c p{2cm} p{6cm}}
    \toprule
    \textbf{Momento} & \textbf{Giorno/Ora} & \textbf{Tipo} & \textbf{Durata} & \textbf{Descrizione} \\
    \midrule
    \endfirsthead

    \toprule
    \textbf{Momento} & \textbf{Giorno/Ora} & \textbf{Tipo} & \textbf{Durata} & \textbf{Descrizione} \\
    \midrule
    \endhead

    \midrule
    \multicolumn{5}{r}{\textit{Continua nella pagina successiva}} \\
    \midrule
    \endfoot

    \bottomrule
    \endlastfoot
    Sprint Retrospective & Martedì, 14:30 & Sincrono & $\sim$ 30min & Riflessione interna del team per individuare punti di forza, criticità e possibili miglioramenti da applicare nello sprint successivo tramite azioni correttive. \\

    Sprint Review & Martedì, 15:00 & Sincrono & $\sim$ 20min & Condivisione dei risultati ottenuti durante lo sprint, verifica di ciò che è stato prodotto rispetto alla Definition of Done e chiusura formale dello sprint precedente.\\

    Sprint Planning & Martedì, 15:20 & Sincrono & $\sim$ 2h & Definizione degli obiettivi dello sprint, aggiornamento del backlog, assegnazione dei task tramite Google Forms e stima dei tempi previsti. \\
    
    Daily Stand-up & Ogni giorno, 14:30 & Asincrono & $\sim$ 2min per persona & Ogni membro aggiorna il canale con: (1) cosa ho fatto, (2) cosa farò, (3) blocchi o dubbi (non obbligatorio nei giorni in cui c'è riunione sincrona). \\
    
\end{longtable}

\subsection{Risorse Umane}
Il team è composto da 6 membri, con ruoli rotativi secondo il modello Scrum. Ogni membro dedica in media dalle 10 alle 20 ore settimanali al progetto.

\subsection{Strumenti e tecnologie}
\begin{itemize}
    \item \textbf{Codice sorgente:} GitHub 
    \item \textbf{Task management:} GitHub Projects + Google Form (creazione e chiusura task) 
    \item \textbf{Comunicazione sincrona:} Riunioni settimanali (martedì + venerdì) 
    \item \textbf{Comunicazione asincrona:} WhatsApp (daily asincrono) 
    \item \textbf{Documentazione:} {\LaTeX} su VS Code con estensione per autocompilazione PDF + repository GitHub 
    \item \textbf{CI/CD:} GitHub Actions 
    \item \textbf{Pianificazione visiva:} \link{https://www.onlinegantt.com/\#/gantt}{OnlineGantt}
    \item \textbf{Tecnologie applicative}: Da definire
\end{itemize}